% chr19.tex -- owner: More RLBWT Tricks task.
% Extracted by the merging task (2026-10-30) from paper.tex: Sections 3-7 (Section 5 shrunk to two
% sentences, Table 2 shortened to its 1% rows), with the unitig example and Figure 1 from the
% introduction.  Renamings applied (see NOTES.md).  Lines marked %MERGE were added or rewritten by the
% merging task; please check them.
\section{Long matches: crossing them quickly}
\label{sec:chr19}

%MERGE
Our third experiment finds the MEMs of whole chromosome~19 haplotypes
with respect to collections of 16 to 64 other chromosome~19s
(\Cref{sec:chr19-exp}). Almost every character of such a pattern lies
in a long MEM, so the matches are easy to find and we do not classify
them; the bottleneck is crossing them, since backward search takes one
step per character even when the pattern matches material that occurs
many times in the collection. In this section $T = S_1 \# S_2 \# \cdots \#
S_d \# \$$ is a collection of $d$ genomes over
$\{\texttt{A},\texttt{C},\texttt{G},\texttt{N},\texttt{T}\}$, with $\$
< \# < \texttt{A} < \texttt{C} < \texttt{G} < \texttt{N} <
\texttt{T}$, on the forward strand only, and we index $T$ and the
reversed collection $T'$ (\Cref{sec:prelim}), each with an RLBWT.

Repetitiveness is not only a challenge for indexes built on RLBWTs but also an
opportunity. One reason the $r$-index was surprising is that it locates
the occurrences of a pattern not only in space proportional to $r$ but
also in predecessor time, $O(\log \log_w (n/r))$ per occurrence on a
word RAM with $w$-bit words~\cite{GagieNavarroPrezza2020}, and Nishimoto
and Tabei~\cite{NishimotoTabei2021} reduced this to constant time per
occurrence, still in $O(r)$ words, and backward search to constant time
per step when the alphabet has polylogarithmic size. FM-indexes and
compressed suffix arrays that use $O(n \log \sigma)$ bits access the
suffix array in $O((\log_\sigma n)^\epsilon)$ time, for a constant
$\epsilon > 0$~\cite{FerraginaManzini2000,GrossiVitter2005}, and Kempa and
Kociumaka~\cite{KempaKociumaka2026} recently proved that, for binary
texts in general, constant-time access requires $\Omega(n \log^\epsilon
n)$ bits for some $\epsilon > 0$. On a text that is not repetitive,
$O(r)$ words are $\Theta(n \log n)$ bits, so the comparison makes sense
only because the text is repetitive. Nishimoto and Tabei obtained their
bounds with move structures, which replace the rank queries of backward
search by lookups in a table with one entry per run: for a text that is
not repetitive the table would be huge, but for a pangenome it is
small, and in practice much faster than rank
queries~\cite{ZakeriEtAl2024}. Because pangenomes compress so well, we
have room to store structures that would make no sense for a single
reference, and the door is open to doing things with RLBWTs that we
could not do with BWTs. This paper takes another step in that
direction.

Ben Langmead asked us whether BWT-based indexes could get the kind of
speedup that unitigs give $k$-mer indexes. \Cref{fig:dbg} shows a collection of twelve strings
and its de Bruijn graph for $k = 3$, uncompacted and compacted. Every
string contains the unitig \texttt{AACGCTC}, and ten of them contain
the string \texttt{GTCAAGG}, which would be a unitig if two strings did
not have \texttt{GTCCAGG} instead; we call such a string an
\emph{almost-unitig}. The unitig and the almost-unitig appear in
different orders in different strings, so neither can be extended.
Every occurrence of a unitig's last $k$-mer is preceded by the rest of
the unitig, unless a string starts inside the unitig; in the BWT this means
that the interval of the last $k$-mer is mapped one-to-one, in order,
onto the interval of the whole unitig, so a backward search that reaches
the last $k$-mer of \texttt{AACGCTC} and then reads \texttt{AACG} in the
pattern can cross those four characters in one step. For an
almost-unitig, some rows drop out on the way, and a bitvector over the
interval of the last $k$-mer says which.

\begin{figure}[t]
  \centering
  \small
  \begin{tabular}{@{}l r@{}}
    \texttt{GTA\,AACGCTC\,GTCAAGG\,TGG} & $\times 5$ \\
    \texttt{GTA\,AACGCTC\,GTCCAGG\,TGG} & $\times 1$ \\
    \texttt{GTA\,GTCAAGG\,AACGCTC\,TGG} & $\times 5$ \\
    \texttt{GTA\,GTCCAGG\,AACGCTC\,TGG} & $\times 1$
  \end{tabular}

  \medskip
  \begin{tikzpicture}[x=0.62cm,y=0.62cm,>={Stealth[length=1.4mm]},v/.style={draw,rounded corners=1pt,inner sep=1pt,font=\tiny\ttfamily}]
\node[v,fill=gray!20] (GTA) at (0.00,0.00) {GTA};
\node[v,fill=gray!20] (AAC) at (3.00,1.60) {AAC};
\node[v,fill=gray!20] (CTC) at (7.00,1.60) {CTC};
\node[v,fill=gray!20] (GTC) at (3.00,-1.60) {GTC};
\node[v,fill=gray!20] (AGG) at (7.00,-1.60) {AGG};
\node[v,fill=gray!20] (TGG) at (10.00,0.00) {TGG};
\node[v] (TAA) at (1.00,0.53) {TAA};
\node[v] (AAA) at (2.00,1.07) {AAA};
\node[v] (TAG) at (1.00,-0.53) {TAG};
\node[v] (AGT) at (2.00,-1.07) {AGT};
\node[v] (ACG) at (4.00,1.60) {ACG};
\node[v] (CGC) at (5.00,1.60) {CGC};
\node[v] (GCT) at (6.00,1.60) {GCT};
\node[v] (TCG) at (5.67,0.53) {TCG};
\node[v] (CGT) at (4.33,-0.53) {CGT};
\node[v] (TCT) at (8.00,1.07) {TCT};
\node[v] (CTG) at (9.00,0.53) {CTG};
\node[v] (TCA) at (4.00,-1.60) {TCA};
\node[v] (CAA) at (5.00,-1.60) {CAA};
\node[v] (AAG) at (6.00,-1.60) {AAG};
\node[v] (TCC) at (4.00,-2.60) {TCC};
\node[v] (CCA) at (5.00,-2.60) {CCA};
\node[v] (CAG) at (6.00,-2.60) {CAG};
\node[v] (GGA) at (5.67,-0.53) {GGA};
\node[v] (GAA) at (4.33,0.53) {GAA};
\node[v] (GGT) at (8.00,-1.07) {GGT};
\node[v] (GTG) at (9.00,-0.53) {GTG};
\draw[->] (GTA) -- (TAA);
\draw[->] (GTA) -- (TAG);
\draw[->] (TAA) -- (AAA);
\draw[->] (AAA) -- (AAC);
\draw[->] (AAC) -- (ACG);
\draw[->] (ACG) -- (CGC);
\draw[->] (CGC) -- (GCT);
\draw[->] (GCT) -- (CTC);
\draw[->] (CTC) -- (TCG);
\draw[->] (CTC) -- (TCT);
\draw[->] (TCG) -- (CGT);
\draw[->] (CGT) -- (GTC);
\draw[->] (GTC) -- (TCA);
\draw[->] (GTC) -- (TCC);
\draw[->] (TCA) -- (CAA);
\draw[->] (CAA) -- (AAG);
\draw[->] (AAG) -- (AGG);
\draw[->] (AGG) -- (GGT);
\draw[->] (AGG) -- (GGA);
\draw[->] (GGT) -- (GTG);
\draw[->] (GTG) -- (TGG);
\draw[->] (TCC) -- (CCA);
\draw[->] (CCA) -- (CAG);
\draw[->] (CAG) -- (AGG);
\draw[->] (TAG) -- (AGT);
\draw[->] (AGT) -- (GTC);
\draw[->] (GGA) -- (GAA);
\draw[->] (GAA) -- (AAC);
\draw[->] (TCT) -- (CTG);
\draw[->] (CTG) -- (TGG);
\end{tikzpicture}


  \medskip
  \begin{tikzpicture}[x=0.62cm,y=0.62cm,>={Stealth[length=1.4mm]},v/.style={draw,rounded corners=1pt,inner sep=1.5pt,font=\scriptsize\ttfamily}]
\node[v] (GTA) at (0,0) {GTA};
\node[v] (TAAA) at (1.5,0.8) {TAAA};
\node[v] (TAGT) at (1.5,-0.8) {TAGT};
\node[v,fill=gray!20] (AACGCTC) at (5,1.6) {AACGCTC};
\node[v] (TCGT) at (5.6,0.45) {TCGT};
\node[v] (GGAA) at (4.4,0.45) {GGAA};
\node[v,fill=gray!8] (GTC) at (3,-1.6) {GTC};
\node[v,fill=gray!8] (TCAAG) at (5,-1.6) {TCAAG};
\node[v,fill=gray!8] (TCCAG) at (5,-2.6) {TCCAG};
\node[v,fill=gray!8] (AGG) at (7,-1.6) {AGG};
\node[v] (TCTG) at (8.5,0.8) {TCTG};
\node[v] (GGTG) at (8.5,-0.8) {GGTG};
\node[v] (TGG) at (10,0) {TGG};
\draw[->] (GTA) -- (TAAA);
\draw[->] (GTA) -- (TAGT);
\draw[->] (TAAA) -- (AACGCTC);
\draw[->] (TAGT) -- (GTC);
\draw[->] (AACGCTC) -- (TCGT);
\draw[->] (AACGCTC) -- (TCTG);
\draw[->] (TCGT) -- (GTC);
\draw[->] (GTC) -- node[font=\tiny,fill=white,inner sep=0.5pt] {10} (TCAAG);
\draw[->] (GTC) -- node[font=\tiny,fill=white,inner sep=0.5pt] {2} (TCCAG);
\draw[->] (TCAAG) -- node[font=\tiny,fill=white,inner sep=0.5pt] {10} (AGG);
\draw[->] (TCCAG) -- node[font=\tiny,fill=white,inner sep=0.5pt] {2} (AGG);
\draw[->] (AGG) -- (GGAA);
\draw[->] (AGG) -- (GGTG);
\draw[->] (GGAA) -- (AACGCTC);
\draw[->] (TCTG) -- (TGG);
\draw[->] (GGTG) -- (TGG);
\end{tikzpicture}

  \caption{Twelve strings (top; spaces added for readability) and their
    de Bruijn graph for $k = 3$, uncompacted (middle; branching
    $k$-mers shaded) and compacted (bottom). \texttt{AACGCTC} is a
    unitig. \texttt{GTCAAGG} is an almost-unitig: ten strings contain
    it, and two have \texttt{GTCCAGG} instead (edge labels count the
    strings).}
  \label{fig:dbg}
\end{figure}

%MERGE: condensed from the introduction of More RLBWT Tricks
We built these \emph{shortcuts} (\Cref{sec:shortcuts}), generalised to
windows of a reference sequence, and used them to speed up
forward--backward search and BML over the RLBWTs of $T$ and $T'$. They
are beaten, however, by \emph{two-level
indexing}~\cite{BrownEtAlRLCSA,HongEtAl2024,DengEtAl2022}, which also
indexes a parse of the collection (\Cref{sec:twolevel}). Our main
contribution here (\Cref{sec:blocks}) is an encoding of the parse into
blocks of phrases that are each the lexicographically smallest
successor of the phrase before them, which lets the search cross a
whole block in one step while all three indexes stay in the same order.

\subsection{Shortcuts}
\label{sec:shortcuts}

The following observation lets a backward search notice that it has
matched a string $\beta$ from a precomputed set.

\begin{lemma}
\label{lem:detect}
If $|\beta| \le |X|$ and $\intv(X)$ is non-empty, then $\intv(X)
\subseteq \intv(\beta)$ if and only if $\beta$ is a prefix of $X$.
\end{lemma}

\begin{proof}
If $\beta$ is a prefix of $X$, every suffix starting with $X$ starts
with $\beta$. Otherwise the prefix $\beta'$ of $X$ of length $|\beta|$
differs from $\beta$, so $\intv(X) \subseteq \intv(\beta')$, which is
disjoint from $\intv(\beta)$.
\end{proof}

Let $\beta$ be the last $k$-mer of a unitig or almost-unitig $\alpha
\beta$. We store $\intv(\beta)$, the first row $f$ of $\intv(\alpha
\beta)$ and a bitvector $M$ over $\intv(\beta)$ marking the rows whose
suffixes are preceded by $\alpha$ (for a unitig, all of them, and $M$
need not be stored). When the current match $X$ has $\intv(X) =
[\mathrm{lo}, \mathrm{hi}] \subseteq \intv(\beta)$ and the pattern has
$\alpha$ immediately to the left of $X$, which we check by comparing
characters, the search jumps to
\[
\intv(\alpha X) = [f + \mathrm{rank}_1(M, \mathrm{lo}),\ f +
\mathrm{rank}_1(M, \mathrm{hi} + 1) - 1],
\]
where ranks are relative to the start of $\intv(\beta)$: by the order
preservation of LF, the marked rows of $\intv(\beta)$ are mapped in
order onto $\intv(\alpha\beta)$, and those inside $\intv(X)$ onto
$\intv(\alpha X)$. If the result is empty, the search takes ordinary
steps. Since the jump is taken only after $\alpha$ has been confirmed in
the pattern, the search returns the same result as plain backward
search, and it can also return the list of shortcuts it took. This list
is a path through the de Bruijn graph, or, with the windows below,
through a pangenome graph whose vertices are the windows, so a BWT-based
index can report the path of a match along with its interval; we find
this of more practical than theoretical interest, and do not develop
it further here.

In a pangenome, the strings that commonly precede a window are
naturally chosen from a reference $\Rref$ in the collection, such as
CHM13~\cite{NurkEtAl2022}. In the example of \Cref{fig:dbg}, with the
first string as $\Rref$ and windows of length $\ell = 4$ starting at
multiples of $4$, a search that has matched the window
\texttt{GTGG} at position $16$ can check whether the pattern has
\texttt{CAAG} to its left and, if so, cross it in one step; of the six
rows of $\intv(\texttt{GTGG})$, the one from the string with
\texttt{GTCCAGG} drops out. We stored, for each window $\beta_p = \Rref[p..p
+ \ell)$ whose interval is disjoint from those of the other windows,
shortcuts across the $2^t \ell$ characters $\Rref[p - 2^t \ell .. p)$
for $0 \le t \le \tau$ with $2^t \ell$ dividing $p$; a shortcut ends at
the start of another window, so shortcuts chain. At a window, the
search tries the longest shortcut whose string the pattern matches, and
then shorter ones.

\begin{lemma}
\label{lem:count}
Suppose $P[j - m\ell .. j + \ell) = \Rref[p - m\ell .. p + \ell)$,
the current match $P[j..|P|)$ has its interval inside that of
$\beta_p$, $P[j - m\ell .. |P|)$ occurs in $T$, and none of the windows
at $p - (m-1)\ell, \ldots, p$ was dropped because it collides with
another. Then the search reaches or passes position $j - m\ell$ of $P$
with at most $2\lceil \log_2(m+1) \rceil + \lfloor m / 2^\tau \rfloor$
shortcuts.
\end{lemma}

The proof is in \Cref{app:proofs}.

The rows of
$\intv(\beta_p)$ preceded by a longer string are a subset of those
preceded by a shorter one, so instead of one bitvector per length we
store, for each row, the number of lengths it survives, in a wavelet
tree; with $\tau = 0$ this is a single bitvector, as for an
almost-unitig. With the RLBWTs of $T$ and $T'$, and windows at the same
boundaries of $\Rref$ in both, forward and backward searches take
shortcuts in both directions.

On the chromosome~19 data of \Cref{sec:chr19-exp}, with $\ell = 25$
and $\tau = 6$, shortcuts reduce the backward steps per pattern
character from about $3.7$ to $0.45$ and make BML with $L = 1000$
$7.8$ times faster, for $1.7$ times the space of the two
RLBWTs; other values of $\ell$ and $\tau$ were slower (shorter windows
collide more often, and longer ones are crossed less often). This
speedup stays about the same as the collection grows from 16 to 64
individuals, but the space grows to $2.5$ times that of the
RLBWTs, mostly for the per-row symbols (\Cref{tab:main}).

\subsection{Two-level indexing with a syncmer parse}
\label{sec:twolevel}

Two-level indexing beats shortcuts on our data, so we describe it in
detail. Fix $k > s$. A $k$-mer over
$\{\texttt{A},\texttt{C},\texttt{G},\texttt{T}\}$ is a \emph{closed
syncmer}~\cite{Edgar2021} if its smallest $s$-mer, under an order $g$ on
$s$-mers, is its first or its last (ties included). We take $g(\gamma)$
to be a hash of the smaller of the 2-bit codes of $\gamma$ and of its
reverse $\gamma^{\mathrm{rev}}$, so $g(\gamma) =
g(\gamma^{\mathrm{rev}})$, and a $k$-mer is a closed syncmer if and only
if its reverse is. We use $k = 51$ and $s = 12$.

\begin{lemma}[Edgar~\cite{Edgar2021}]
\label{lem:window}
Among any $k - s$ consecutive $k$-mers of a string over
$\{\texttt{A},\texttt{C},\texttt{G},\texttt{T}\}$ there is a closed
syncmer.
\end{lemma}

\begin{proof}
The $k$-mers starting at positions $0, \ldots, k - s - 1$ of a string
$Y$ of length $2k - s - 1$ contain the $2k - 2s$ $s$-mers starting at
positions $0, \ldots, 2k - 2s - 1$. Let $q$ be the start of a smallest
of them. If $q \le k - s - 1$, it is the first $s$-mer of the $k$-mer
starting at $q$; otherwise it is the last $s$-mer of the $k$-mer
starting at $q + s - k \in [0, k - s - 1]$. In both cases it is
smallest in that $k$-mer.
\end{proof}

\paragraph*{The parse.}
Let $T^*$ be $T$ or $T'$. The \emph{boundaries} of $T^*$ are $b_0 = 0$ and
the starting positions of the closed syncmers in $T^*$, $b_0 < b_1 <
\cdots < b_{m-1}$. The phrases are $T^*[b_i .. b_{i+1} + k)$ for $i < m -
1$ and $T^*[b_{m-1} .. n)$, so consecutive phrases overlap by $k$
characters. We number the distinct phrases $1, \ldots, \sigma$ in the
lexicographic order of their strings, and the \emph{parse} $\Pa$ of
$T^*$ is the sequence of their ids, followed by a sentinel $0$. Within a
stretch of $T^*$ over $\{\texttt{A},\texttt{C},\texttt{G},\texttt{T}\}$,
\Cref{lem:window} bounds the length of a phrase by $2k - s$. By
the symmetry of $g$, the boundaries of $T'$ mirror those of $T$, and
the phrases of $T'$ are the reverses of those of $T$, except at the
ends; so one dictionary serves both directions.

Eddie Ferro noted that the parse at closed syncmers is
prefix-free~\cite{FerroPC}; we include a proof for completeness.

\begin{lemma}
\label{lem:order}
No phrase is a proper prefix of another, and for boundaries $b_i \ne
b_j$, $T^*[b_i .. n) < T^*[b_j .. n)$ if and only if $\Pa[i .. m] <
\Pa[j .. m]$.
\end{lemma}

\begin{proof}
A phrase contains no syncmer start strictly between its first character
and its last $k$ characters, while the last $k$-mer of every phrase
other than the last is a closed syncmer, and whether a $k$-mer is a
closed syncmer does not depend on its context. So if a phrase other than the last were a proper prefix of another
phrase, the last $k$-mer of the shorter phrase would be a syncmer
starting strictly inside the longer one;
and the last phrase ends with \$, which occurs nowhere else. Now let $h$
be the first index with $\Pa[i + h] \ne \Pa[j + h]$, which exists since
the sentinel occurs once. The phrases at $i, \ldots, i + h - 1$ and $j,
\ldots, j + h - 1$ are equal, so $T^*[b_i .. b_{i+h} + k) = T^*[b_j ..
b_{j+h} + k)$ and $b_{i+h} - b_i = b_{j+h} - b_j$. The phrases at $i +
h$ and $j + h$ are different and neither is a prefix of the other, so
they differ at some offset smaller than both lengths, and $T^*[b_i..n)$
and $T^*[b_j..n)$ compare as these two phrases do.
\end{proof}

\paragraph*{Mapping between the two levels.}
Let $B$ be the bitvector over the rows of the RLBWT of $T^*$ with $B[i] =
1$ if and only if $\mathrm{SA}[i]$ is a boundary. Every occurrence of a
closed syncmer is a boundary, so $B$ is the union of the intervals of
the distinct closed syncmers in $T^*$ plus the row of position $0$, and it
has at most that many runs of 1s; we store it run-length coded. By
\Cref{lem:order}, the $u$th 1 in $B$ (counting from $0$) is the row of
the parse with rank $u + 1$, row $0$ being that of the sentinel.

\begin{lemma}
\label{lem:bounce}
Let $q < q'$ be consecutive starts of closed syncmers in a pattern
$P[0..e)$, with $q' + k \le e$. The occurrences of $P[q..e)$ in $T^*$ are
exactly the boundaries $b_t$ such that the phrase at $t$ is $P[q .. q'
+ k)$ and $P[q'..e)$ occurs at $b_{t+1}$. Moreover, for any string $Y$
that starts with a closed syncmer, the rows of $\intv(Y)$ are all
marked in $B$, and the corresponding rows of the parse are
$[\mathrm{rank}_1(B, \mathrm{lo}) + 1, \mathrm{rank}_1(B, \mathrm{hi} +
1)]$ for $\intv(Y) = [\mathrm{lo}, \mathrm{hi}]$.
\end{lemma}

\begin{proof}
An occurrence of $P[q..e)$ at position $a$ of $T^*$ starts with a closed
syncmer, so $a$ is a boundary $b_t$. The $k$-mers of $T^*$ starting at
$a + 1, \ldots, a + q' - q$ lie inside the occurrence, so they are
closed syncmers exactly when the corresponding $k$-mers of $P$ are:
only the last is. Hence $b_{t+1} = a + q' - q$ and the phrase at $t$ is
$P[q..q' + k)$. The converse is immediate. The second claim holds
because every occurrence of $Y$ starts at a boundary, and by
\Cref{lem:order} the order of rows marked in $B$ is that of the parse.
\end{proof}

\paragraph*{Searching.}
To find the longest suffix of $P[x..e)$ that occurs in $T^*$, and its
interval, we find the starts of the closed syncmers in $P$ and look up
the ids of the phrases $P[q .. q' + k)$ between consecutive starts in a
dictionary of the phrases of $T$; the same lookup gives the id of the
reversed phrase in $T'$, so the pattern is parsed once for both
directions. The search takes backward steps on the RLBWT from $e$ to
the last syncmer start $q$ with $q + k \le e$, maps the interval to the
parse by \Cref{lem:bounce}, takes one backward step on the BWT of the
parse for each phrase of the pattern to the left, and, when a phrase is
not in the dictionary or a step on the parse fails, maps back with
select queries on $B$ and continues with backward steps on the RLBWT.

\begin{theorem}
\label{thm:twolevel}
If the dictionary lookups are correct, this search returns the same
longest suffix and interval as backward search. If $P[x..e)$ is over
$\{\texttt{A},\texttt{C},\texttt{G},\texttt{T}\}$, it takes at most $3k
- 2s - 1$ backward steps on the RLBWT, plus one step on the parse per
phrase of the pattern crossed and one that may fail.
\end{theorem}

\begin{proof}
Correctness follows from \Cref{lem:bounce} by induction over the
phrases. By \Cref{lem:window}, if $e - x \ge 2k - s - 1$ then one of the
last $k - s$ $k$-mers of $P[x..e)$ is a closed syncmer, so the search
reaches the parse after at most $2k - s - 1$ steps; otherwise it takes
at most $2k - s - 2$ steps in total. At the left end, either the search
is within $k - s - 1$ characters of $x$, by \Cref{lem:window}, or the
phrase $P[q .. q' + k)$ to the left of the current syncmer start $q'$
does not precede it in any occurrence, in which case $P[q..e)$ does not
occur by \Cref{lem:bounce} and the backward steps from $q'$ fail before
reaching $q$, after fewer than $q' - q \le k - s$ steps.
\end{proof}

Each step on the parse, like each step on the block sequence of
\Cref{sec:blocks}, is a predecessor query in one Elias--Fano list of
the runs of its BWT, grouped by symbol (\Cref{sec:prelim}); with
$r_\Pa$ runs in the BWT of the parse, it takes $O(1 + \log(\sigma m /
r_\Pa))$ time, and it touches only the runs of one symbol.

The dictionary stores the 64-bit hashes of the phrases of $T$, sorted
and split into buckets by their leading bits, with the phrases' ids in
$T$ and $T'$; it checks the whole hash, and we check at construction
time that the phrases of $T$ have distinct hashes. A phrase of the
pattern that is not a phrase of $T$ is then mistaken for one only if
its hash collides with one of $\sigma$ stored values, so if the hash
function behaves like a random function, the lookups are correct with
high probability; we have never observed an error. On our data
this dictionary takes $46$\,MiB for both directions, against
$96$\,MiB per direction for a hash table with the phrase ids, and
looking up all the phrases of the patterns, with prefetching, takes
about $1.1$ seconds, which our timings include. Splitting the patterns at runs of \texttt{N}
matters for the bound: the current parse does not split phrases at
\texttt{N}, so a phrase can be as long as a run of \texttt{N} in the
collection.

\subsection{Blocks}
\label{sec:blocks}

\paragraph*{Encodings of the parse.}
\label{sec:encodings}
%MERGE: More RLBWT Tricks' Section 5, shrunk to two sentences as agreed
After two-level indexing a search takes one step per phrase, and most
phrases of a pattern are predictable from the phrase before them, as
the characters of a unitig are. We first tried replacing each phrase by
its rank by frequency among the phrases that follow the phrase before
it~\cite{GagieHybrid2026}, but then the encoded suffixes are not in the
order of the suffixes of the parse, so the search cannot bounce between
the levels with $B$; the blocks below keep the first phrase of each
block explicit, which keeps all the indexes in the same order.

For a phrase id $x$, let $\succp(x)$ be the smallest id $y$ such that
$x y$ occurs in $\Pa$. Position $i > 0$ of $\Pa$ is a
\emph{continuation} if $\Pa[i] = \succp(\Pa[i-1])$. A \emph{block} is a
maximal substring $\Pa[a .. a + j]$ such that $a$ is not a continuation
and $a + 1, \ldots, a + j$ are; since $\succp$ is a function, the block
is determined by its \emph{symbol} $(\Pa[a], j)$, and is $\delta, \succp(\delta),
\ldots, \succp^j(\delta)$ for $\delta = \Pa[a]$. The \emph{block sequence}
$\Delta$ lists the symbols of the blocks of $\Pa$ in order, followed by
a sentinel. We order symbols by $\delta$ ascending and then by $j$
descending.

\begin{lemma}
\label{lem:blockorder}
For block starts $a \ne a'$ of $\Pa$, $\Pa[a..m] < \Pa[a'..m]$ if and
only if the suffix of $\Delta$ starting at the block of $a$ is smaller
than the one starting at the block of $a'$.
\end{lemma}

\begin{proof}
Equal symbols are equal strings of phrases followed by block starts,
so consider the first blocks that differ, $(\delta, j)$ at $a$ and $(\delta', j')$
at $a'$. If $\delta \ne \delta'$, both orders are decided by $\delta$ and $\delta'$.
Otherwise, by symmetry, let $j > j'$. The two suffixes agree on $j' + 1$
phrases, and the first continues with $\succp^{j'+1}(\delta)$. The block at
$a'$ is not the last block of $\Pa$, since the last phrase of $\Pa$
occurs only once and would then also be in the block at $a$; so the
second suffix continues with the first phrase $y$ of its next block,
which follows $\succp^{j'}(\delta)$ in $\Pa$ but is not a continuation, so $y
> \succp^{j'+1}(\delta)$ by the definition of $\succp$. So the first suffix is smaller, as its symbol is.
\end{proof}

So if $BS$ marks the rows of the parse whose suffixes start at block
starts, including the row of the sentinel, the $u$th marked row
(counting from $0$) corresponds to row $u$ of the BWT of $\Delta$, and a \emph{block step} on that BWT, with a symbol $(\delta, j)$,
prepends a whole block to the parse suffixes at block starts. We store
$\succp$, $BS$ (run-length coded when that is smaller, which on our
data it is by a factor of $9$ to $20$), the BWT of $\Delta$, the
symbols grouped by $\delta$ with their $j$'s in decreasing order, and one bit
per phrase saying whether it is a continuation anywhere in $\Pa$.

\begin{lemma}
\label{lem:space}
Let $\eta$ be the number of distinct closed syncmers in $T^*$, $\sigma$
the number of distinct phrases, $m$ the number of phrases, $r_\Pa$ and
$r_\Delta$ the numbers of runs in the BWTs of $\Pa$ and $\Delta$, $b$
the number of blocks and $s'$ the number of distinct block symbols.
Besides the RLBWT of $T^*$ and the dictionary, the parse index of $T^*$,
with blocks, takes
\[
O\bigl(\eta \log(n/\eta) + r_\Pa \log(\sigma m / r_\Pa) + r_\Delta
\log(s' b / r_\Delta) + b \log(m/b) + \sigma \log \sigma + s' \log m\bigr)
\]
bits. The dictionary, shared by $T$ and $T'$, takes $\sigma (64 - h + 2
\lceil \log_2 (\sigma + 3) \rceil) + O(\sigma)$ bits, where $2^h$, the
number of its buckets, is about $\sigma / 8$.
\end{lemma}

The proof is in \Cref{app:proofs}.

Nothing here depends on $n$ except through $\eta$, and $\eta$,
$\sigma$ and the numbers of runs grow slowly as haplotypes are added
(\Cref{sec:chr19-exp}). With 16 individuals, the parse indexes of $T$
and $T'$ take $102.5$\,MiB without blocks and $139.6$\,MiB with them,
of which the dictionary takes $46$\,MiB.

A pattern is split into blocks in the same way, using $\succp$ of the
collection. The key observation is that a block boundary in the
pattern is a block boundary in every occurrence that extends one
phrase to its left.

\begin{lemma}
\label{lem:blockstart}
If phrase $i$ of the pattern is not a continuation of phrase $i - 1$,
then in every occurrence in $\Pa$ of the pattern's phrases $i - 1, i,
\ldots$, the position of phrase $i$ is a block start.
\end{lemma}

\begin{proof}
The phrase before it in $\Pa$ is the pattern's phrase $i - 1$, and the
pattern's phrase $i$ is not $\succp$ of it.
\end{proof}

The search with block jumps proceeds as in \Cref{sec:twolevel}, except
that when the current parse interval is exact for the pattern's
phrases from $i$ on, and $i$ is a block start of the pattern whose
phrase is complete, it restricts the interval to block starts with two
rank queries on $BS$ and takes block steps, one per block of the
pattern, for as long as they succeed. In an occurrence of the pattern's
phrases from $i$ on that starts at a block start of $\Pa$, the block
boundaries of $\Pa$ are those of the pattern, since $\succp$ is a
function and by \Cref{lem:blockstart}; so, by induction, each interval
after a block step consists exactly of the occurrences of the pattern's
phrases from the current position on that start at block starts of
$\Pa$, and by \Cref{lem:blockstart} it contains every occurrence that
extends one more phrase to the left. The interval after the last block
step, however, may miss occurrences whose first phrase is a
continuation in $\Pa$. We fix this with the following rule.

\begin{lemma}
\label{lem:backup}
Suppose a sequence of $h \ge 1$ block steps ends at the pattern's
phrase $i$, the last of them crossing phrases $i, \ldots, i'$. If the
phrase $i$ of the pattern is never a continuation in $\Pa$, the parse
rows between the first and the last row of the final interval form
the exact parse interval for phrases $i, i+1, \ldots$. Otherwise,
taking parse steps for phrases $i', \ldots, i$ from the parse rows
between the first and the last row of the interval before the last
block step gives exact intervals.
\end{lemma}

\begin{proof}
In the first case, every occurrence starts at a block start, so the
final interval contains all of them, and the parse rows between its
first and last rows are all occurrences, by contiguity. In the second
case, let $J$ be the interval before the last block step, for the
pattern's phrases from $i' + 1$ on; it contains every occurrence
preceded by phrase $i'$, and the parse rows between its first and last
rows are occurrences of phrases $i' + 1, \ldots$. The rows of the true
interval outside that range are not preceded by phrase $i'$, by
\Cref{lem:blockstart}, so the parse step for $i'$ is exact, and so are
the following parse steps.
\end{proof}

After backing up, the search takes no block steps until it has crossed
phrase $i$ again. Each phrase is crossed by at most one block step and
one parse step, and each sequence of block steps ends with at most one
failed step, so \Cref{thm:twolevel} holds with at most two steps per
phrase on the parse and on $\Delta$. On our data, with 16 individuals, block jumps reduce the steps on
the parse from $23.3$ million to $3.1$ million, plus $2.5$ million block
steps, for BML with $L = 1000$.

\paragraph*{Other continuations.}
Blocks play the role of unitigs in the parse, except that a block
follows the smallest successor of each phrase rather than its only one.
The smallest successor is not the only choice that keeps the indexes in
the same order. Let $\chi$ assign each phrase $x$ one of the phrases that
follow it in $\Pa$, call position $i$ a continuation if $\Pa[i] =
\chi(\Pa[i-1])$, and give each symbol a third component, a bit saying
whether the first phrase of the next block is smaller or larger than $\chi$
of the block's last phrase. If symbols with the same $\delta$ are ordered
with the smaller ones first, by $j$ ascending, and then the larger
ones, by $j$ descending, the proof of \Cref{lem:blockorder} goes
through, and \Cref{lem:blockstart,lem:backup} do not depend on the
choice of $\chi$. With the smallest successor every bit says ``larger'', so
no bits are stored. Choosing the most frequent successor maximises the
number of continuations, and its blocks may shrink less as haplotypes
are added, since the most frequent successor of a phrase changes less
than its smallest one. We have implemented this choice and checked that
the search remains exact. On chromosome~19 with 16 individuals, it makes
blocks $3.3$ times longer, $31.4$ phrases on average instead of $9.5$,
and its blocks hardly shrink as haplotypes are added: $30.9$ phrases
with 32 individuals and $30.5$ with 64, against $8.3$ and $7.1$ with
the smallest successor. It does not yet pay off, however: a search also
takes parse steps for up to about a block at each end and when it backs
up, so with 16 individuals it takes $7.7$ million parse steps and $0.6$
million block steps for BML with $L = 1000$, instead of $3.1$ and $2.5$
million, and it is $5$ to $6\%$ slower.

\subsection{Experiments}
\label{sec:chr19-exp}

\paragraph*{Data.}
From the assemblies of release~2 of the Human Pangenome Reference
Consortium~\cite{HPRC2025,LucasEtAl2026} we took chromosome~19. The
collections consist of the reference CHM13~\cite{NurkEtAl2022}, of
$61.7$ million bases, and both haplotypes of the first 16, 32 or 64
individuals, by sample identifier, whose haplotypes each have
chromosome~19 as a single sequence of at least 90\% of the length of
CHM13; with 16 individuals there are 33 sequences and $n =
2{,}032{,}742{,}766$ characters. The patterns are haplotype~1 of the
next two such individuals, HG00423 and HG00438, of $61.7$ and $61.0$
million bases, split at runs of \texttt{N} into $4$ patterns of
$122.4$ million characters in all. All
sequences were already in the orientation of CHM13.

\paragraph*{Implementation and setup.}
The wavelet tree of our RLBWT uses rank9-style
bitvectors~\cite{Vigna2008}. We
built the RLBWTs with Big-BWT~\cite{BoucherEtAl2019}, and the parse
indexes from the collection and the RLBWTs in a streaming pass.
%MERGE: code citation moved here from the conclusions
Our code for this section is available at
\url{https://github.com/TravisGagie/MEM-finding}, in the directory \texttt{chr19}. All experiments ran
on a desktop with an Intel Core i7-8700 CPU (3.2\,GHz, 6 cores) and 16\,GiB of RAM, under Ubuntu 26.04, with
the searches pinned to one core and a single thread for every tool; we
report the median of five runs. All the methods we compare could be run
in parallel in the same way, by searching different patterns, or
different parts of a pattern, in different threads. Every configuration of ours gave the
same MEMs, with the same intervals, as plain search. For ropebwt3, which
indexes both strands and reports a MEM if it or its reverse complement
matches, we checked that all differences are reverse-strand matches:
with 16 individuals and $L = 100$, for example, ropebwt3 reports 234
MEMs that occur only on the reverse strand, and our 140 MEMs that
ropebwt3 does not report lie inside them, and with $L = 1000$ the
numbers are 5 and 2. Movi~2~\cite{ZakeriEtAl2026}, which we ran with
its regular index of both strands and a table of 10-mers, reads
\texttt{N} in the collection as \texttt{A}, and most of its differences
from our MEMs are explained in the same way, but not all: with
$L = 100$, at each size, between 140 and 281 strings reported only by
Movi~2 and 16 MEMs reported only by us remain, and with $L = 1000$ at
most 89 and 1. Many of the former are not maximal; for example, they
start a few characters after a MEM that has the same end. We have not
investigated these differences.
Ropebwt3 and Movi~2 answer a slightly more general question than our
index of $T$ and $T'$, which finds only forward-strand MEMs. Building
ropebwt3's index took $20$ and $43$ minutes with 32 and 64 individuals,
with a peak of $9.6$\,GiB.
Given the RLBWTs, building the parse indexes for both directions took
$5.4$, $7.5$ and $12.1$ minutes for 16, 32 and 64 individuals, with a
peak memory use of $1.9$, $3.5$ and $7.0$\,GiB.

\begin{table}[t]
  \centering
  \caption{Finding the MEMs of two chromosome~19 haplotypes (122.4
    million characters) with respect to collections of 16, 32 and 64
    individuals: space in MiB for both directions, and time in seconds,
    single-threaded, for BML with $L = 1000$ and $L = 100$ and for
    forward--backward (FB) with $L = 1000$. Plain search uses only the
    two RLBWTs, the shortcuts have $\ell = 25$ and $\tau = 6$, and
    ropebwt3 and Movi~2 index both strands. For Movi~2, the FB column
    gives the time to find all MEMs.}
  \label{tab:main}
  \begin{tabular}{l l r rrr}
    \toprule
    individuals & method & MiB & BML, 1000 & BML, 100 & FB, 1000 \\
    \midrule
    16 & plain search            & 314 & 290.5 & 287.3 & 287.2 \\
       & shortcuts               & 543 & 37.2 & 39.5 & 39.3 \\
       & two-level               & 417 & 12.3 & 14.5 & 13.8 \\
       & two-level, block jumps  & 454 & 8.0 & 10.3 & 9.7 \\
       & ropebwt3                & 157 & 64.3 & 65.5 & 63.1 \\
       & Movi~2                  & 554 & 25.9 & 28.9 & 57.3 \\
    \midrule
    32 & plain search            & 349 & 321.5 & 318.2 & 318.6 \\
       & shortcuts               & 699 & 40.3 & 42.2 & 41.5 \\
       & two-level               & 475 & 12.1 & 13.7 & 12.8 \\
       & two-level, block jumps  & 526 & 7.1 & 8.7 & 7.9 \\
       & ropebwt3                & 178 & 72.8 & 73.9 & 69.0 \\
       & Movi~2                  & 586 & 30.2 & 34.2 & 65.1 \\
    \midrule
    64 & plain search            & 406 & 348.7 & 345.5 & 344.3 \\
       & shortcuts               & 998 & 43.1 & 44.8 & 44.2 \\
       & two-level               & 569 & 12.9 & 14.3 & 13.4 \\
       & two-level, block jumps  & 642 & 7.1 & 8.6 & 7.7 \\
       & ropebwt3                & 212 & 79.5 & 80.2 & 76.8 \\
       & Movi~2                  & 639 & 35.5 & 40.8 & 85.1 \\
    \bottomrule
  \end{tabular}
\end{table}

\paragraph*{Results.}
\Cref{tab:main} shows the space and time of each method. With 16
individuals, the parse has $106.9$ million phrases, $4.17$ million of
them distinct, and $11.2$ million blocks, so a block has $9.5$ phrases
on average. Without shortcuts, finding the MEMs with BML and $L = 1000$
takes $3.7$ backward steps per pattern character, since a forward search
and two backward searches cover each MEM and consecutive MEMs overlap.
Two-level indexing replaces most of them by $0.19$ steps on the parse
per character, block jumps reduce these to $0.05$ steps on the parse and
on $\Delta$, and only $0.06$ backward steps per character remain on the
RLBWT. Two-level indexing with block jumps is then $36$ times faster
than plain search and $4.6$ times faster than the shortcuts, for $1.4$
times the space of the RLBWTs, less than the shortcuts take. It is $6.4$
to $8.0$ times faster than ropebwt3, whose index of both strands takes a
third of its space, and $2.8$ to $3.2$ times faster than Movi~2, whose
index of both strands, with its table of 10-mers, takes $1.2$ times its
space. With 32 and 64 individuals, it is $8.5$ to $10.3$ and $9.3$ to
$11.2$ times faster than ropebwt3, whose index takes a third of its
space, and $3.9$ to $4.3$ and $4.7$ to $5.0$ times faster than Movi~2,
whose index takes $1.1$ and $1.0$ times its space.

KeBaB~\cite{BrownEtAl2025KeBaB}, with $k = 20$, does not help here,
as it does with long reads that mostly do not match the collection: its
fragments keep more than $99.5\%$ of the characters with 16 individuals,
and more than $99.8\%$ with 32 and 64, since almost every
character of a whole chromosome lies in a long MEM, and neither plain
search, nor block jumps, nor ropebwt3 gets faster on the fragments. We
checked that the MEMs found in the fragments are exactly the MEMs of
length at least $L$.

As the collection grows from 16 to 64 individuals, the average number of
phrases per block drops from $9.5$ to $8.3$ and $7.1$: each new
haplotype adds successors to some phrases, and a phrase that used to be
the only successor of its predecessor may no longer be the smallest one.
Nevertheless, block jumps get faster: for BML with $L = 1000$ their
speedup over plain search grows from $36$ to $45$ and $49$, and over
two-level indexing without them from $1.5$ to $1.7$ and $1.8$. The
number of steps on the parse and on $\Delta$ stays at about $5$ million,
with more block steps and fewer parse steps, while plain search and
two-level indexing without jumps take more steps, presumably because the
MEMs get longer and overlap more. The space of the parse index grows
much more slowly than that of the shortcuts: from $0.44$ to $0.51$ and
$0.58$ times that of the RLBWTs, against $0.73$, $1.00$ and $1.46$.
